\item Dos pulsos A y B se mueven en direcciones opuestas a lo largo de una cuerda tensa con una rapidez $v = 2.00\text{ cm/s}$. La amplitud de A es el doble de la amplitud de B. Los pulsos se muestran en la figura P17.3 en $t = 0$. Bosqueje la onda resultante en $t = 1.00\text{ s}$, $1.50\text{ s}$, $2.00\text{ s}$, $2.50\text{ s}$ y $3.00\text{ s}$.